\section{Conclusions}
\label{sec:conclusions}

We have addressed the problem of choosing the size of a permissioned BFT
validator set for an e-governance workload in the realistic setting where no
large testbed is available. \emph{First}, the identifiability theorem shows that
a single-factor scaling exponent converges to $\gamma\lambda-\beta$, where
$\lambda$ is a property of the measurement design rather than of the protocol; a
positive exponent is therefore consistent with a load generator scaled alongside
the network and is not by itself evidence that throughput grows with the
validator count. This offers a concrete explanation for the discrepancy between
the predictions of our earlier work~\cite{peniak2026a,peniak2026b} and indicates
which of the two designs was affected. \emph{Second}, the resource-normalized
measurement protocol fixes the per-validator compute budget independently of the
validator count, separating consensus overhead from host contention and making
the required factorial design --- and hence the theorem --- usable on a single
commodity machine; where exact $\q$-invariance is rejected, the fitted law is
reported at the reference quota. \emph{Third}, calibrated lognormal prediction
intervals,
validated by empirical coverage on held-out validator counts, replace point
accuracy figures with a falsifiable statement. \emph{Fourth}, coupling the
throughput law with a Hoeffding-type detection bound turns sizing into a
chance-constrained program whose feasible set is a closed interval
$[\nmin,\nmax]$ computable in $O(\log\n)$ evaluations, with absolute demand
entering through an explicit production scale-up $\kappaScale$; the detection
bound itself is standard, and what it contributes is a constraint monotone in the
direction opposite to throughput, so that an empty interval certifies a
structural rather than a provisioning failure.

Of the four, we regard the first as the one most likely to outlive the specific
measurements reported here, since it is a statement about estimation rather than
about any particular system. The entire campaign executes in a public
continuous-integration workflow on one four-core runner, so every result can be
reproduced without dedicated infrastructure.

\paragraph{Future work.}
Protocol-level instrumentation covering equivocation faults; validation across
several consensus implementations; extension of the feasibility window to
latency and finality constraints; and online recalibration of $\beta$ as a
deployment grows.

\subsection*{Conflict of interest}
The authors declare no conflict of interest.

\subsection*{Funding}
The research was conducted without financial support.

\subsection*{Data availability}
The workflow definition, network and load generators, fault-injection scripts,
detector, analysis code and raw measurement artefacts are available in the
project repository.

\subsection*{Use of artificial intelligence}
AI-assisted tools were used for language editing and typesetting. All scientific
results, proofs, experimental design and analysis were produced by the authors.

\subsection*{Contributions of authors}
\textbf{Peniak B.\,O.}: conceptualization, methodology, formal analysis,
software, investigation, data curation, writing --- original draft,
visualization.
\textbf{Liubinskyj B.\,B.}: supervision, validation, writing --- review and
editing.
